\section{Related work}
\label{sec:related}

\paragraph{Interference from updated values.} Proactive interference grows with the number of overwrites of
a key, and errors are mostly earlier values of the same key \citep{wang2025unable}. Models also retrieve
the first value of a key more reliably than the latest \citep{chattaraj2026dual,qiao2026retrieval}.
\citet{guo2026context} name the failure \emph{stale binding}. They find that almost all errors on current
value queries return an old value of the queried variable, that the current value stays linearly
recoverable on failures, that moving a historical mention near the query increases errors and historical
wording reduces them, and that test-time attention routing repairs most errors in 3--9B models. Their
decision scenarios, run with four closed reasoning models, showed no difference between full histories
and resolved snapshots. \citet{qian2026forget} and \citet{chao2026stale} evaluate forgetting and stale
memories in dialogue and agent settings. STALE in particular separates recognizing an update from applying
it in a later policy. Both are accuracy benchmarks with naturalistic text. Our setting differs in three
ways: a single update instead of long overwrite chains, decisions whose correct action is fixed by an
explicit policy, and matched minimal pairs with lexical controls. The failures we study sit below the
interference loads at which retrieval breaks.

\paragraph{Entrainment and distraction.} Tokens that appear in the context gain probability whether or
not they are relevant, random tokens included \citep{niu2025entrainment}. Irrelevant context also degrades
reasoning \citep{shi2023distracted}. Our entity-mention and unattached-word constructions measure this
effect in the update setting. The question we ask is whether a superseded assignment adds decision
pressure beyond it. \citet{yang2026compared} argue that counterfactual prompt edits are compound
treatments and must be compared against matched alternative edits rather than against no edit. Our design
follows that logic: each construction applies the same lexical edit, and the contrasts compare constructions.

\paragraph{Binding and state tracking.} Mechanistic work identifies how models bind attributes to entities
\citep{feng2024bind,gurarieh2026mixing,prakash2026lookbacks}. It also shows that models aggregate state
changes at query time rather than updating incrementally \citep{tang2026entities,kim2023entity}. These
accounts offer plausible reasons why lexical identity and position could matter more than assignment
history, but our experiments are behavioral and do not test them.

\paragraph{Knowledge conflicts.} Our records contain an intra-context conflict with an explicit temporal
resolution \citep{xu2024conflicts}. Related benchmarks ask for the value valid at a given time
\citep{yang2026conflicts}. That is our historical retrieval check, not our decision task.
